\documentstyle[%


righttag%


,11pt%


,amssymb%


,russian%               remove in Eng.


,izvru%                 Set izven% in Eng.


]{amsart}





% 


%% !! {,} repl. for . in eng.





\newcommand{\tts}[1]{}


\numberwithin{equation}{section}





\theoremstyle{plain}


\theorembodyfont{\it}


\newtheorem{thms}{\hskip\parindent Теорема}[section]








%\hoffset=-15mm%                  For Eng.


\setcounter{page}{3}


\god{2003}


\nomer{1~(488)} %                        Remove in Eng.


%\nomer{Vol.\,47, No.\,1}%               Set in Eng.


%\str{\pageref{begin}--\pageref{end}}%  Set in Eng.


\udk{519.63}





\author{\it В.Н.\,Абрашин, Н.Г.\,Жадаева}


% NB:   ^^ remove in Eng.


\title{Об аддитивных итерационных методах \\


и оценках их скорости сходимости}





\date{}


%\pagestyle{myheadings}%         Set in Eng.


\pagestyle{plain} %              Remove in Eng.


\begin{document}


%\thispagestyle{plain}%          Set in Eng.


\maketitle


\label{begin}











\section*{Введение}





Многочисленные итерационные методы позволяют дать явный 


вид приближения к решению стационарных задач при помощи эволюционных систем   


уравнений. Это дает возможность получить быстро сходящиеся итерационные


алгоритмы с довольно простой их реализацией.





Среди эффективных классических итерационных процессов широкое применение имеет 


метод переменных направлений, который основан на специальных подходах 


релаксации исходной задачи с возможной редукцией сложной задачи к 


последовательности


простейших [1]--[5]. Все эти методы можно отнести к методам расщепления.


Метод переменных направлений можно 


отнести к методам полной аппроксимации, т.\,к.


его алгоритм аппроксимирует исходное уравнение. К его недостатку можно отнести 


ограничение на количество операторов 


расщепления (их должно быть не больше двух).


К методам расщепления, которые можно эффективно использовать в качестве 


итерационных для решения стационарных задач, следует отнести


методы факторизации [6]--[8] и стабилизирующей поправки [9], [10]. 


При многокомпонентном расщеплении эти методы требуют попарной коммутируемости 


пространственных операторов [11], [12]. 


Возможность применения методов расщепления


(дробных шагов) [1], [4] как итерационных методов решения стационарных задач


без требования коммутируемости операторов расщепления была показана в 


[13], однако вопросы скорости сходимости не были изучены. В данной работе 


изучены аддитивные итерационные методы полной аппроксимации для некоммутируемых 


операторов расщепления [14], [15]. 


Предложенные алгоритмы развивают известные методы 


расщепления, для них получены оценки скорости сходимости и показано их 


преимущество перед классическими. 





\section{Постановка задачи и аддитивные итерационные методы}





Рассмотрим операторное уравнение первого рода


\begin{equation}


Au=f


\end{equation}


с линейным оператором $A : H \rightarrow H$ (не обязательно дискретным),


действующим в вещественном гильбертовом пространстве $H$ со скалярным


произведением $(u,v)$ и нормой $\|u\|=\sqrt{(u,u)}$.


Пусть $A$ --- положительно


определенный оператор, $A\ge cE$, $c>0$, $H_A$ --- пространство $H$,


снабженное скалярным произведением $(u,v)_A=(Au,v)$ 


и нормой $\|u\|_A=\sqrt{(Au,u)}$.





Решение многих стационарных задач (1.1) с положительным оператором можно


рассматривать как предельный переход при $t\rightarrow \infty$


нестационарной эволюционной задачи


\begin{equation}


\frac {dy}{dt}+Ay=f,\quad t>0,\quad y(0)=v_0.


\end{equation}


При решении стационарной задачи с помощью эволюционной задачи (1.2)


промежуточные значения решения этого уравнения не представляют интереса.


Интересуют лишь асимптотические решения 


задачи (1.2) при $t\rightarrow \infty$.


По сути дела в этом и состоит единство и различие этих задач [13].





Решения задач (1.1), (1.2) связаны неравенством вида


\begin{equation}


\|y(t)-u\|\le e^{-ct} \|v_0-u\|.


\end{equation}


Эта схема позволяет конструировать эффективные итерационные методы, 


в том числе и на основе неявных методов. Например, для 


решения эволюционной


задачи (1.2) можно использовать неявный разностный метод


\begin{equation}


(\wh y-y)/\tau+A\wh y=f,\quad y(0)=y_0,


\end{equation}


на сетке $t=t_j=j\tau$, $y=y(t)$, $\wh y=y(t+\tau)$, $j\ge 0$.


Разностная схема асимптотически устойчива при любых $\tau>0$. Разностная


задача (1.4), как и (1.2), может служить 


алгоритмом для решения уравнения (1.1),


при этом естественно нет необходимости говорить о точности (1.4) при решении


задачи (1.2), особенно при больших $t$.





Для разностной эволюционной задачи (1.4) имеет место неравенство типа (1.3),


которое имеет вид


\begin{equation}


\|y(t_j)-u\|\le e^{-\delta t_j} \|y_0-u\|,\quad \delta=\delta(\tau)>0.


\end{equation}


Величина $\delta(\tau)$ при больших $t$ может существенно отличаться от


$c$, однако для решения задачи (1.1) 


это не важно. Использование (1.4) в качестве


итерационного метода практически невыгодно, 


т.\,к. при решении этого уравнения 


количество операций на каждом шаге итерационного процесса не меньше, чем


при непосредственном решении задачи (1.1). 


Поэтому представляет интерес построение


таких итерационных методов, в которых сохранялась 


бы скорость сходимости, близкая к


(1.5), и количество итераций на 


каждом шаге итерационного метода существенно


уменьшалось. Такими методами являются аддитивные разностные схемы, 


предложенные в [16], [17].


Аддитивные методы базируются на представлении оператора $A$ в виде


$\sum\limits_{\alpha=1} ^p A_\alpha$, где каждый из операторов $A_\alpha$


является стационарным (возможно вырожденным) 


оператором. Уравнение (1.1) можно записать 


в виде $\sum\limits_{\alpha=1}^p A_\alpha u_\alpha =f,$ (если $u_\alpha=u$,  


$\alpha=\overline {1,p}$, то это равенство переходит в уравнение (1.1)).


Вместо скалярной эволюционной задачи рассмотрим 


задачу Коши для системы уравнений


\begin{equation}


\frac{d{y_\alpha}}{dt}+\sum\limits_{\beta=1}^p A_\beta y_\beta =f,\quad


\alpha=\overline {1,p},\quad y_\alpha(0)=y_0.


\end{equation}


При решении (1.6) можно применять различные разностные схемы и, если они


являются асимптотически устойчивыми, их можно 


использовать в качестве итерационных


методов решения уравнения (1.1). В 


данной работе детально изучается алгоритм параллельного


вычисления, предложенный в [18], [19], который имеет вид


\begin{equation}


\Big(\overset{s+1}{y}_{\!\!\!\alpha}-\overset{s}{\wt y}{}\Big)/\tau+


\sigma A_\alpha\Big(\overset{s+1}{y}_{\!\!\!\alpha}


-\overset{s}{y}{}_\alpha\Big)+


\sum\limits_{\beta=1}^p A_\beta \overset{s}{y}{}_\beta=f,\quad


\overset{s}{\wt y}{}=


p^{-1}\sum\limits_{\beta=1}^p \overset{s}{y}{}_\beta,\quad


\alpha=\overline {1,p},\quad\overset{0}{y}{}_\alpha=y_0.


\end{equation}


Для погрешности данного итерационного 


метода $\overset{s}{z}{}_\alpha=u-\overset{s}{y}{}_\alpha$


имеет место задача 


\begin{equation} 


\Big(\overset{s+1}{z}{}_{\!\!\!\alpha}-\overset{s}{\wt z}{}\Big)+ 


\sigma A_\alpha\Big(\overset{s+1}{z}{}_{\!\!\!\alpha}


-\overset{s}{z}{}_\alpha\Big)+


\sum\limits_{\beta=1}^p A_\beta \overset{s}{z}{}_\beta=0,\quad


\overset{s}{\wt z}{}=


u-p^{-1}\sum\limits_{\beta=1}^p \overset{s}{z}{}_\beta,\quad


\alpha=\overline {1,p},\quad\overset{0}{z}{}_\alpha=u-y_0.


\end{equation}


Отметим, что итерационный метод (1.7) изучался в [17]--[22].


Исследование этого метода на базе уравнения (1.8) 


позволяет получить эффективные оценки погрешности


метода и доказать их высокую скорость сходимости и фактически получить оценки 


погрешности, близкие к (1.5).





Ниже будут рассмотрены случаи коммутируемых операторов $A_\alpha$,


некоммутируемых операторов $A_\alpha$, а 


также вырожденных операторов расщепления.


Будут также приведены примеры применения операторных итерационных


методов (1.7) к решению стационарных задач математической физики.





\section{Оценка скорости сходимости аддитивного итерационного метода}





Вопросам исследования экономичных итерационных методов посвящена 


обширная литература [2]--[4]. Большое внимание уделено повышению скорости


сходимости этих методов за счет выбора 


оптимального итерационного параметра $\tau$.


Причем надо отметить, что коммутативный случай изучен более детально. Это


касается прежде всего двухкомпонентного метода переменных направлений и


метода факторизации. В [21], [22] изучены вопросы сходимости для 


многокомпонентного аддитивного итерационного метода полной аппроксимации вида 


\begin{equation}


\Big(\overset{s+1}{y}_{\!\!\!\alpha}


-\overset{s}{y}{}_\alpha\Big)/\tau+


\sigma A_\alpha\Big(\overset{s+1}{y}_{\!\!\!\alpha}


-\overset{s}{y}{}_\alpha\Big)+


\sum\limits_{\beta=1}^p A_\beta \overset{s}{y}{}_\beta=f,\quad


\alpha=\overline {1,p},\quad\overset{0}{y}{}_\alpha=y_0,


\end{equation}


однако получить оценки скорости сходимости 


не удалось. Незначительные изменения алгоритма (2.1), а именно,


усреднение решения в первом выражении, привели к итерационному методу (1.7),


что существенно изменило ситуацию. 


Дело в том, что в итерационном методе компоненты 


решения сходятся друг к другу довольно быстро, а к единому решению ---


медленно. Предложенная в [18]--[20] простая модификация (2.1) вида (1.7)


исправила эту ситуацию. В данной 


работе изучен вопрос о скорости сходимости итерационного


метода (2.1) и получены оценки погрешности 


невязки в несколько нетрадиционной норме.


Ниже приведены оценки скорости сходимости итерационного процесса (1.7) в 


естественной (легко проверяемой) норме.





Пусть $A_\alpha \ge{\sigma}_\alpha E$, $\sigma_\alpha>0$,


$\alpha=\overline {1,p}$, --- положительно 


определенные и попарно коммутируемые 


операторы, $A_\alpha A_\beta=A_\beta A_\alpha$.


Рассмотрим подробно случай $\sigma=p$. Умножим уравнение (1.8) скалярно


на $\tau A_\alpha \overset{s+1}{z}_{\!\!\!\alpha}$ и просуммируем


по $\alpha=\overline {1,p}$


\begin{multline}


\sum\limits_{\alpha=1}^p \| \overset {s+1}{z}_{\!\!\!\alpha}\|_


{A_\alpha}^2


-\bigg( \overset {s}{\wt z},\ \sum\limits_{\alpha=1}^p 


A_ \alpha\overset {s+1}{z}_{\!\!\!\alpha}\bigg)+


0.5\tau^3 p\sum\limits_{\alpha=1}^p


\|(A_\alpha \overset {s}{z}_\alpha)_t\|^2


+0.5\tau p \sum\limits_{\alpha=1}^p \|


A_\alpha\overset {s+1}{z}_{\!\!\!\alpha}\|^2+ \\


%


+0.25\tau\bigg(\bigg\| \sum\limits_{\alpha=1}^p


A_\alpha(\overset {s+1}{z}_{\!\!\!\alpha}+


\overset {s}{z}_\alpha)\bigg\|^2 -


\tau^2 \bigg\|\sum\limits_{\alpha=1}^p (A_\alpha


\overset {s}{z}{}_\alpha)_t\bigg\|^2 \bigg)=


0.5\tau p \sum\limits_{\alpha=1}^p 


\|A_\alpha \overset {s}{z}{}_\alpha\|^2,


\end{multline}


где $\overset {s}{z}{}_{\alpha t}=(\overset {s+1}{z}_{\!\!\!\alpha}-


\overset {s}{z}{}_\alpha)/\tau$.


Так как $(\overset {s+1}{\widetilde z}-\overset {s}


{ \widetilde z})/\tau =- \sum\limits_{\alpha=1}^p  A_\alpha


\overset {s+1}{z}_{\!\!\!\alpha}$, то


\begin{gather*}


\bigg(\overset {s}{ \widetilde z}, \sum\limits_{\alpha=1}^p 


A_\alpha\overset {s+1}{z}_{\!\!\!\alpha}\bigg)=


-\frac{1}{\tau}\Big((\overset {s+1}{\widetilde z}-\overset {s}


{\widetilde z}),\overset {s}{\widetilde z}\Big)=


- 0.5\tau^{-1} \Big(\|\overset {s+1}{\widetilde z}\|^2 


- \|\overset {s}{\widetilde z} \|^2 -


\tau^2 \|\overset {s}{\widetilde z_t} \|^2 \Big), \\


%


\overset {s}{\widetilde z}{}_t  = - \sum\limits_{\alpha=1}^p 


A_\alpha \overset {s+1}{z}_{\!\!\!\alpha}=


- 0.5 \bigg(\sum\limits_{\alpha=1}^p


A_\alpha(\overset{s+1}{z}_{\!\!\!\alpha}+


\overset{s}{z}{}_\alpha) + \tau\sum\limits_{\alpha=1}^p (  A_\alpha


\overset{s}{z}{}_\alpha)_t \bigg).


\end{gather*}


Учитывая вышесказанное, из равенства (2.2) получим


\begin{multline}


\tau\sum\limits_{\alpha=1}^p \|


\overset {s+1}{z}_{\!\!\!\alpha}\|^2_{A_\alpha}+ 


0.5\|\overset {s+1}{\widetilde z}\|^2 + 0.5\tau^4


\bigg(p \sum\limits_{\alpha=1}^p \|(A_\alpha


\overset{s}{z}{}_\alpha)_t\|^2- \bigg\|\sum\limits_{\alpha=1}^p (A_\alpha 


\overset{s}{z}{}_\alpha)_t\bigg\|^2 \bigg)+  \\


%


+0.5\tau^2 p\sum\limits_{\alpha=1}^p \|A_\alpha


\overset {s+1}{z}_{\!\!\!\alpha}\|^2+


0.25\tau^2\bigg(\bigg\|\sum\limits_{\alpha=1}^p A_\alpha 


(\overset{s+1}{z}_{\!\!\!\alpha} +\overset{s}{z}{}_\alpha)\bigg\|^2 + 


\tau^2\bigg\|\sum\limits_{\alpha=1}^p


(A_\alpha\overset{s}{z}{}_\alpha)_t\bigg\|^2\bigg)- \\


%


-0.5\tau^2\bigg(\bigg\|\sum_{\alpha=1}^p A_\alpha


\overset{s+1}{z}_{\!\!\!\alpha}


\bigg\|^2-\bigg\|\sum\limits_{\alpha=1}^p


A_\alpha\overset{s}{z}{}_\alpha\bigg\|^2\bigg) = 


0.5\|\overset {s}{\widetilde z} \|^2 +0.5\tau^2 p \sum\limits_{\alpha=1}^p 


\|A_\alpha \overset{s}{z}{}_\alpha\|^2.


\end{multline}


Учитывая, что


\begin{gather*}


(A_\alpha \overset{s}{z}{}_\alpha)_t - (A_\beta \overset{s}{z}{}_\beta)_t = 


-p^{-1} \tau^{-2}\overset{s+1}{v}{}^{(\alpha,\beta)},\quad 


\overset{s+1}{v}{}^{(\alpha,\beta)} =\overset{s+1}{z}_{\!\!\!\alpha} -


\overset{s+1}{z}_{\!\!\beta}, \\


%


0.5p \sum\limits_{\alpha=1}^p \|(A_\alpha\overset{s}{z}{}_\alpha)_t\|^2- 


0.5\bigg\|\sum\limits_{\alpha=1}^p


(A_\alpha \overset{s}{z}{}_\alpha)_t\bigg\|^2 =


\sum\limits_{\alpha,\beta=1,\alpha > \beta}^p


\|(A_\alpha \overset{s}{z}{}_\alpha)_t


-(A_\beta \overset{s}{z}{}_\beta)_t\|^2,


\end{gather*}


будем иметь


\begin{multline}


2\tau\sum\limits_{\alpha=1}^p \|\overset {s+1}{z}{}_{\!\!\!\alpha}


\|^2_{A_\alpha}+ 


\|\overset {s+1}{\widetilde z}\|^2 + p^{-2} \|\overset {s+1}{v} \|^2_3 +


\tau^2 p \sum\limits_{\alpha=1}^p \| A_\alpha


\overset{s+1}{z}_{\!\!\!\alpha}\|^2 -


\tau^2\bigg\|\sum\limits_{\alpha=1}^p


A_\alpha\overset{s+1}{z}_{\!\!\!\alpha}\bigg\|^2 = \\


%


=\|\overset {s}{\widetilde z}\|^2 +\tau^2 p \sum\limits_{\alpha=1}^p 


\|A_\alpha\overset {s}{z}{}_\alpha\|^2,


\quad \|\overset {s+1}{v}\|^2_3 =


\sum\limits_{\alpha,\beta=1,\alpha>\beta}^p


(\overset {s+1}{v}{}^{(\alpha,\beta)},


\overset {s+1}{v}{}^{(\alpha,\beta)}).


\end{multline}


Пусть $\delta=\min\limits_{1<\alpha<p}\delta_\alpha$, 


тогда $2\tau\sum\limits_{\alpha=1}^p 


\|\overset {s+1}{z}_{\!\!\!\alpha}\|^2_{A_\alpha}\ge


2\tau\delta\sum\limits_{\alpha=1}^p 


\|\overset {s+1}{z}_{\!\!\!\alpha}\|^2\ge


2\tau\delta p\|\overset {s}{\widetilde z}\|^2$.





Также имеет место неравенство


$\|A_\alpha z_\alpha\|^2\le \Delta_\alpha (A_\alpha z_\alpha,


z_\alpha)$, где $\Delta_\alpha>A_\alpha$ --- верхняя граница спектра оператора


$A_\alpha$. И если $\Delta=\max\limits_{1<\alpha<p}\Delta_\alpha$, то


$$


\sum\limits_{\alpha=1}^p(A_\alpha z_\alpha,z_\alpha)\ge


\Delta^{-1}\sum\limits_{\alpha=1}^p


\|A_\alpha z_\alpha\|^2.


$$


Из равенства (2.4) следует неравенство вида


\begin{multline}


(1+2\Theta\tau\delta p)\|\overset {s+1}{\widetilde z}\|^2+(1+(1-\Theta)


\Delta^{-1}\tau^{-1})\tau^2p\sum\limits_{\alpha=1}^p\|A_\alpha


\overset {s+1}{z}_{\!\!\!\alpha}\|^2\le\|\overset {s}{\widetilde z}\|^2+ \\


%


+\tau^2p^2\sum\limits_{\alpha=1}^p\|A_\alpha


\overset {s}{z}{}_\alpha\|^2,\quad 0<\Theta<1.


\end{multline}


Отсюда следует 





\begin{thms}


Пусть $A_\alpha\ge\delta_\alpha E$,


$\alpha=\overline{1,p}$, $\delta_\alpha>0$.


Тогда итерационный метод  $(1.7)$


сходится и для его скорости сходимости справедлива оценка


$$


Q(\overset {s+1}{ z})\le q^{s+1} Q(\overset {0}{ z}),\quad s=0, 1,\dotsc,


$$


где


$$


Q(\overset {s+1}{ z})=\|\overset {s+1}{\widetilde z}\|^2+\tau^2


\sum\limits_{\alpha=1}^p\|A_\alpha


\overset {s}{z}{}_\alpha\|^2,\quad q=\min\Big((1+2\tau\delta p)^{-1},


(1+\Delta^{-1}p^{-2}\tau^{-1})\Big).


$$


\end{thms}





Оптимальная скорость сходимости в данном случае достигается при выполнении 


равенства $1+2\tau\theta\delta p=1+(1-\theta)\Delta^{-1}\tau^{-1}$, 


т.\,е. при $\tau=\big(\frac{1-\theta}{2\theta\delta\Delta p}\big)^{1/2}$.


Этот результат согласуется с обычным методом переменных направлений и 


методом факторизации. Однако в данном случае нет ограничения 


на количество операторов расщепления и не требуется их коммутируемость.





Приведем оценку погрешности, полученную для итерационного метода (1.7), и 


сравним ее с оценкой (2.5).





\begin{thms}[{\series{m}\shape{n}\selectfont [21]}]


Пусть $A_\alpha\ge\delta_\alpha E$, $\delta_\alpha\ge\delta>0$,


$\alpha=\overline {1,p}$. Тогда для итерационного метода $(1.7)$ имеет место


\begin{equation}


Q(\overset {s}{y})\le q^{-s}Q(\overset {0}{y}),


\end{equation}


где $Q(\overset {s}{y})=\|r(\overset {s}{y})\|^2+p^{-2}\tau^{-2}


\|\overset {s}{v}\|_3^2$, $q=1+2\delta p\tau$, $r(\overset {s}{y})=


\sum\limits_{\alpha=1}^p A_\alpha \overset {s}{ y_\alpha}-f$.


\end{thms}





Из оценки (2.6), вообще говоря, 


не следует сходимость итерационного метода (1.7)


к решению исходной задачи, т.\,к. невязка метода $r(\overset {s}{y})$


не согласована с естественной невязкой $A \overset{s}{y}-f$, которая 


используется для классических итерационных методов. 


Из оценки (2.6) можно получить оценку, подобную (2.5). Имеет место тождество


$$


\overset {s}{y}{}_\alpha=\overset {s}{\wt y}+p^{-1}\sum\limits_{\beta=1}^p 


(\overset {s}{y}{}_\beta-y_\beta)=\overset {s}{\wt y}+p^{-1}


\sum\limits_{\beta=1}^p \overset {s}{v}{}^{(\alpha,\beta)},


$$


и для $\overset {s}{z}=y-\overset {s}{\wt y}$ получим равенство вида


\begin{equation}


A\overset {s}{z}=-\sum\limits_{\alpha=1}^p A_\alpha \bigg(p^{-1} 


\sum\limits_{\beta=1}^p \overset {s}{v}{}^{(\alpha,\beta)}\bigg)+


r(\overset {s}{y}).


\end{equation}


Из равенства (2.7) находим $\overset {s}{\rho}=-\sum\limits_{\alpha=1}^p


B_\alpha \Big(p^{-1}\sum\limits_{\beta=1}^p


\overset {s}{v}{}^{(\alpha,\beta)}\Big)+


A^{-1} r(\overset {s}{y})$,


где $B_\alpha=A^{-1}A_\alpha=\linebreak


\Big(E+\sum\limits_{\beta=1,\ \alpha\ne\beta}^p


A_\alpha^{-1}A_\beta\Big)$. Здесь возникает необходимость выделить отдельно 


коммутируемый случай. При условии попарной коммутируемости операторов


$A_\alpha$, $\alpha=\overline{1,p}$, имеет место неравенство


$\|B_\alpha\|<1$ и $\|\overset {s}{z}\|\le c^{-1}\|r(\overset {s}{y})\|+


\|\overset {s}{v}\|_3$, $0<c<A$. Из оценки (2.5) следуют неравенства


\begin{equation}


\|\overset {s}{v}\|_3\le p\tau(q^{-s}Q(\overset {0}{y})^{1/2}),\quad


\|r(\overset {s}{y})\| \le Q(\overset {0}{y})^{1/2}.


\end{equation}


Отсюда для $\overset {s}{z}$ имеет место оценка $\| \overset {s}{\rho}\| \le


(p\tau+c^{-1})q^{-s/2}(Q(\overset {0}{y}))^{1/2}.$ 


Для сходимости итерационного метода (1.7) справедлива





\begin{thms}


Если операторы $A_\alpha \ge \delta_\alpha E$,


$\delta_\alpha > \delta$,


$\delta >0$, попарно коммутируемы, то итерационный метод $(1.7)$ при


$\sigma=p$ сходится к решению исходной задачи, 


и для скорости его сходимости справедлива 


оценка


$$


\|\overset {s}{z}\|\le (p\tau+c^{-1})(1+2\sigma p \tau)^{s/2}\|r


(\overset {0}{y})\|.


$$


\end{thms}





Таким образом, в данном случае итерационный метод (1.7) сходится при любых 


$\tau>0$. Оптимальное же значение 


итерационного параметра $\tau$, как и следует из (2.8),


достигается при $\tau=\tau_0=p^{-1}c^{-1}$. Оценки (2.8) показывают, что 


сходимость итерационного метода (1.7) 


зависит только от нижней границы спектра


операторов $A, A_\alpha$. Это говорит о 


том, что скорость сходимости алгоритма (1.7) 


близка к чисто неявной схеме (1.4). 


Операторы $A$, $A_\alpha$, $\alpha=\overline


{1,p}$, могут быть как дискретными, так и непрерывными. В случае разделяющих 


переменных метод (1.7) по своим качествам близок к методу Фурье.





Для оценки числа итераций $S$, необходимого для достижения требуемой


точности, достаточно потребовать, чтобы


$(p\tau+c^{-1})(1+2\delta c p)^{-S/2}<\varepsilon$.


Пусть $c+p\sigma_\alpha,$ тогда при $\tau=p^{-1}c^{-1}$ имеем


\begin{equation}


\rho \ge \rho_0(\varepsilon)=2(\ln (2c^{-1}/\varepsilon))/\ln(1+2p^{-1}).


\end{equation}


Из (2.9) видно, что при дискретных операторах


$A$, $A_\alpha$, $\alpha=\overline {1,p}$, количество 


итераций не зависит от шага пространственной сетки. Это говорит


о преимуществе метода (1.7) перед классическим методом переменных направлений 


и освобождает нас от сложной процедуры выбора 


оптимального итерационного параметра.





В некоммутируемом случае такого результата получить пока не удалось.


В некоммутируемом случае [21] справедлива





\begin{thms}


Если $E \Delta_\alpha \ge A_\alpha \ge \delta_\alpha E$,


$\Delta_\alpha>0$, $\delta_\alpha>0$, то итерационный метод $(1.7)$ сходится к 


решению исходной задачи, и для скорости его сходимости справедлива оценка


\begin{equation}


\|\overset {s}{z}\|_A \le c^{-1/2}\|r(\overset {s}{y})\|+


\bigg(\sum\limits_{\alpha=1}^p \bigg\|p^{-1}


\sum\limits_{\beta=1}^p \overset {s+1}{v}{}^{(\alpha,\beta)}


\bigg\|_A ^2\bigg)^{1/2}


\le (c^{-1/2}+\tau\Delta^{-1/2})q^{s/2}(Q(\overset {0}{y}))^{1/2},


\end{equation}


где $\Delta=\sum\limits_{\alpha=1}^p \Delta_\alpha$.


\end{thms}





Оценка (2.10) типична для метода переменных направлений и метода 


факторизации, но она справедлива в некоммутируемом случае и любом количестве 


слагаемых разбиения оператора~$A$. Это выражается в том, что в (2.10) входит


слагаемое $\tau\Delta^{-1/2}$. 





\section{Некоторые примеры}





Рассмотрим первую краевую задачу для эллиптического уравнения


\begin{equation}


Lu=-\sum\limits_{\alpha,\beta=1}^{p}\frac{\partial}{\partial x_\alpha}


\bigg(k_{\alpha\beta}(x)\,\frac{\partial u}{\partial x_\beta}\bigg)=f(x),


\quad x\in G^{(p)},\quad u(x)=0,\quad x\in\partial G^{(p)},


\end{equation}


где $G^{(p)}=(x_1,\dotsc,x_p)$ --- односвязная область с достаточно гладкой


границей $\partial G^{(p)}$, $L$ --- эл\-лип\-тический оператор, для которого


выполняются условия


\begin{equation}


c_1\sum\limits_{\alpha=1}^{p}\xi_\alpha^2\le


\sum\limits_{\alpha,\beta=1}^{p}k_{\alpha\beta}(x)\


\xi_\alpha\xi_\beta\le c_2\sum\limits_{\alpha=1}^{p}\xi_\alpha^2,


\quad 0\ne\xi=(\xi_1,\dotsc,\xi_p)\in G^{(p)}.


\end{equation}


Предполагается также, что известные функции $f(x)$, $k_{\alpha\beta}(x)$


таковы, что решение задачи (3.1) существует, единственно и обладает


гладкостью, необходимой для корректности рассматриваемых ниже методов.





При $k_{\alpha\beta}=0$, $\alpha\ne\beta$, для решения уравнения (1.1)


можно непосредственно использовать итерационный метод (1.4).


Представление $A=\sum\limits_{\alpha=1}^pA_\alpha$, $A_\alpha\ge0$, не


позволяет применить предложенные методы для уравнений со смешанными


производными. Построим подобное представление в случае, когда в


уравнении (1.1) задан эллиптический дифференциальный оператор со смешанными


производными.





Если $k_{\alpha\beta}(x)\ne k_{\beta\alpha}(x)$, то уравнение (3.1) можно


видоизменить с помощью простой процедуры. Преобразуем оператор $L$,


представив его в виде суммы симметричного и кососимметричного операторов


$L=A+B$, где


\begin{gather*}


Au=\sum\limits_{\alpha,\beta=1}^pA_{\alpha\beta}u, \quad 


A_{\alpha\beta}u=-\frac{\partial}{\partial x_\alpha}


\bigg(a_{\alpha\beta}(x)\,\frac{\partial u}{\partial x_\beta}\bigg), \\


%


a_{\alpha\alpha}=k_{\alpha\alpha}, \quad


a_{\alpha\beta}=0{,}5(k_{\alpha\beta}+k_{\beta\alpha})=a_{\beta\alpha},


\quad \beta\neq\alpha,


\quad \alpha=\overline{1,p},\ \ \beta=\overline{1,p}, \\


%


Bu=\sum\limits_{\alpha=1}^pB_\alpha u, \quad 


B_\alpha u=-0{,}5\bigg(\frac{\partial}{\partial x_\alpha}


\,(b_\alpha(x)u)+b_\alpha(x)\,\frac{\partial u}{\partial x_\alpha}\bigg),\\


%


b_\alpha(x)=0{,}5\sum\limits_{\beta=1,\,\alpha\ne\beta}^p


\frac{\partial}{\partial x_\alpha}\,(k_{\beta\alpha}-k_{\alpha\beta}),


\quad \Vec b(x)=(b_1,\dotsc,b_p), \quad \mathrm{div}\,\Vec b(x)=0


\quad \forall\,x\in G^{(p)}.


\end{gather*}


При этом в пространстве $H$ операторы $A$, $B$ 


обладают следующими свойствами:


1)~$A$ --- самосопряженный положительно-определенный 


оператор, т.\,е. $A=A^*>0$; 


2)~$B=-B^*$, $B_\alpha=-B_\alpha^*$, $\alpha=\overline{1,p}$, и


$(B_\alpha u,u)=0$ для всех $u\in H$.





Далее оператор $A$ разобьем на сумму двух операторов: $A=A^++A^-$ так, что


\begin{gather*}


A^\pm=\sum\limits_{\alpha=1}^{p} A_\alpha^\pm, \quad 


A_\alpha^-u=\sum\limits_{i=1}^{\alpha}A^-_{\alpha i}u, \quad 


A_\alpha^+u=\sum\limits_{i=\alpha}^{p}A^+_{\alpha i}u, \\


%


A_{\alpha i}=\begin{cases}


A_{\alpha i}^+, & i\ge\alpha; \\ A_{\alpha i}^-, & i\le\alpha, \end{cases}


\quad


A_{\alpha i}^\pm u=\frac{\partial}{\partial x_\alpha}


\bigg(a_{\alpha i}^\pm\,\frac{\partial u}{\partial x_i}\bigg),


\quad a_{\alpha\alpha}^-=a_{\alpha\alpha}^+=0{,}5a_{\alpha\alpha}, \quad


a_{\alpha i}=\begin{cases}


a_{\alpha i}^+, & i>\alpha; \\ a_{\alpha i}^-, & i<\alpha. \end{cases}


\end{gather*}





Пусть $H^{2p}$ --- гильбертово пространство вектор-функций


$U=(u_1^\to,\dotsc,u_p^\to,u_p^\gets,\dotsc,u_1^\gets)$,


$u_\alpha^\to,u_\alpha^\gets\in H$, $\alpha=\overline{1,p}$,


обращающихся в нуль на границе $\partial G^{(p)}$, со скалярным


произведением $(U,V)=\sum\limits_{\alpha=1}^{2p}(u_\alpha^*,v_\alpha^*)$


и нормой $\|U\|=\sqrt{(U,V)}$, где $u_\alpha^*=u_\alpha^\to$,


$\alpha=\overline{1,p}$, и $u_\alpha^*=u_{2p+1-\alpha}^\gets$,


$\alpha=\overline{p+1,2p}$.





После преобразований оператора $L$ вместо задачи (3.1) рассмотрим систему


дифференциальных уравнений


\begin{equation}


\frac{\partial y_\alpha}{\partial t}+\sum\limits_{\beta=1}^{2p}


(A_\beta^*y_{(\beta)}+\tfrac{1}{2}B_\beta y_\beta)=f(x),


\ \ x\in G^{(p)},\ \ y_\alpha(0)=y_0,\ \ y_\alpha(x)=0,


\ \ x\in\partial G^{(p)},\ \  \alpha=\overline{1,2p}.


\end{equation}


В системе (3.3) приняты следующие обозначения:


\begin{equation}


A_\beta^*y_{(\beta)}=\begin{cases}


A_\beta^-y_{(\beta)}^\to, & \beta=\overline{1,p}; \\


A_{2p+1-\beta}^+y_{(2p+1-\beta)}^\gets, & \beta=\overline{p+1,2p},


\end{cases} \quad


B_\beta y_{(\beta)}=\begin{cases}


B_\beta y_{(\beta)}^\to, & \beta=\overline{1,p}; \\


B_{2p+1-\beta}y_{(2p+1-\beta)}^\gets, & \beta=\overline{p+1,2p},


\end{cases}


\end{equation}


где $A_\beta^-y_{(\beta)}^\to=A_{\beta\beta}^-y_\beta^\to+


\sum\limits_{i=1}^{\beta-1}A_{\beta i}y_i^\to$,


$A_\beta^+y_{(\beta)}^\gets=A_{\beta\beta}^+y_\beta^\gets+


\sum\limits_{i=\beta+1}^pA_{\beta i}y_i^\gets$.





Можно показать [15], что для компонент произвольной вектор-функции (3.4)


в пространстве $H^{2p}$ справедливо неравенство


$$


\sum\limits_{\beta=1}^{2p}(A_\beta^*y_{(\beta)},y_\beta)\ge


0{,}5c_1\sum\limits_{\beta=1}^{2p}\|\partial y_\beta/\partial x_\beta\|^2.


$$


Аналогичное представление имеет место и для дискретного аналога уравнения


(3.1) [14]. Оператор $A_h$ в данном случае представим в аддитивном виде


$$


A=\sum\limits_{\alpha=1}^p(A_{\alpha h}^-+B_{\alpha h})+


\sum\limits_{\alpha=p+1}^{2p}(A_{2p+1-\alpha


,h}^++B_{2p+1-\alpha,h}).


$$


Для уравнения (3.2) аналог экономичной разностной схемы (1.7) будет  


иметь вид


\begin{gather}


(\widehat y{}_{\beta }^\to-\widehat {\widetilde y})/\tau+


\sigma\tau (A_{\beta h}^- y_{(\beta)}^\to)_t+\sum\limits_{\alpha,i=1}^p


(A_{\alpha h}^-y_{(\alpha)}^\to+A_{i h}^+y_{(i)}^\gets)=0,\notag \\


%


(\widehat  y{}_{\gamma }^\gets-\widehat{\widetilde y})/\tau+


\sigma\tau (A_{\gamma h}^+ y_{(\gamma)}^\gets)_t+\sum\limits_{\alpha,i=1}^p


(A_{\alpha h}^-y_{(\alpha)}^\to+A_{i h}^+y_{(i)}^\gets)=0, \\


%


\beta=\overline{1,p}, \quad \gamma=p+1-\beta, \quad\widehat{\widetilde y}=


(2p)^{-1}\bigg(\sum\limits_{\alpha=1}^p


y_{\alpha}^\to+\sum\limits_{\alpha=1+p}^{2p}


y_{(\alpha)}^\gets \bigg).\notag


\end{gather}


Эти уравнения решаются независимо, каждое из них последовательно. Вместо (3.5)


можно также использовать алгоритм вида


\begin{equation}


(\widehat y{}_{\beta }^\to-\widehat {\widetilde y})/\tau+


\sigma\tau (A_{\beta h}^- y_{(\beta)}^\to)_t+\sum\limits_{\alpha=1}^{2p}


A_{\alpha h}^-y_{(\alpha)}=0,


\end{equation}


$A_{\alpha i h}=A_{\alpha i h}^+$, $i\ge\alpha$,


$A_{\alpha i h}=A_{\alpha i h}^-$, $i\le\alpha$.


Для сходимости итерационных методов (3.5), (3.6) справедлива оценка,


подобная оценке в теореме 2.1.





Проиллюстрируем вышесказанное на примере трехмерной


задачи Дирихле для уравнения Пуассона в единичном кубе


\begin{equation}


\sum_{\alpha=1}^3\frac{\partial^2u}{\partial x_\alpha^2}=


-f(x),\quad x\in G,\quad u(x)=g(x),\quad x\in\Gamma.


\end{equation}


Для приближенного решения соответствующей разностной задачи Дирихле 


использовался трехкомпонентный итерационный


метод переменных направлений (1.7).


Расчеты проводились при заданной точности $\varepsilon=10^{-4}$.


Итерации регулировались условием 


$\|{\overset{s}{\rho}}\|\leq c^{-1}\|r(s)\|+


\|{\overset{s}{v}}\|_3^{~}\leq \varepsilon$.


Вычисления показали, что при малых значениях параметра $\tau$


норма разности компонент метода $\|{\overset{s}{v}}\|_3$ мала, но


невязка $\|r(s)\|$ достаточно велика вследствие относительно большого


значения коэффициента сжатия $1/q$. При больших значениях $\tau$


наблюдается обратная ситуация, т.\,е. норма невязки мала ($1/q$ мало),


а разность компонент велика. Минимальное число итераций достигается


для значения $\tau=\tau_0\sim 0{,}025$, при этом 


$s_0(\varepsilon)\sim 13$. С увеличением количества точек


дискретизации $N$ число итераций, необходимых для достижения заданной


точности $\varepsilon$, не возрастает в соответствии с оценкой .


Таким образом, скорость сходимости итерационного


ММПН (1.7) зависит только от нижней границы спектра оператора $A$ и 


заданной точности и не зависит от шага пространственной сетки.


Если сравнивать метод (1.7) с классическим методом переменных направлений


в двумерном варианте или методом факторизации для случая коммутируемых


операторов разбиения, то скорость сходимости этих методов


зависит от шага пространственной дискретизации


даже при оптимальном выборе итерационных параметров.





Рассмотрим теперь применение метода (1.7) для разностного аналога


уравнения (3.7) в случае некоммутируемых операторов $A_\alpha$ (напр.,


для уравнения с переменными коэффициентами). Как было отмечено,


теорема 2.4 позволяет эффективно вычислять в процессе итераций норму


погрешности метода


$\|{\overset{s}{\rho}}\|_A=\|{\overset{s}{\widetilde y}}


-y\|_{W_2^1}$. Условие $\|{\overset{s}{\rho}}\|_A\leq \varepsilon$


использовалось в качестве критерия для окончания процесса итераций.


Расчеты проводились на квадратной сетке для $\varepsilon=h^2$. Результаты


вычислений занесены в таблицу.


\medskip





\begin{center}


\begin{tabular}{|c|c|c|c|c|c|c|c|}


\hline


\multicolumn{4}{|c|}{$h=0{,}1$}


&


\multicolumn{4}{|c|}{$h=0{,}03$} \\


\hline


$\tau$ & $s$ & $\|r(s)\|$ & $\|v\|_3$ &


$\tau$ & $s$ & $\|r(s)\|$ & $\|v\|_3$ \\


\hline


0{,}003 & 85 & $4{,}9\cdot 10^{-2}$ & $1{,}1\cdot 10^{-8}$ &


0,003 & 115 & $3{,}7\cdot 10^{-3}$ & $2{,}3\cdot 10^{-7}$ \\


\hline


0{,}008 & 35 & $4{,}0\cdot 10^{-2}$ & $3{,}5\cdot 10^{-6}$ &


0{,}005 & 75 & $2{,}1\cdot 10^{-3}$ & $4{,}4\cdot 10^{-5}$ \\


\hline


0{,}02 & 20 & $7{,}5\cdot 10^{-3}$ & $3{,}0\cdot 10^{-4}$ &


0{,}008 & 69 & $6{,}5\cdot 10^{-5}$ & $1{,}1\cdot 10^{-5}$ \\


\hline


0{,}1 & 33 & $6{,}3\cdot 10^{-4}$ & $5{,}7\cdot 10^{-4}$ &


0{,}02 & 94 & $1{,}7\cdot 10^{-5}$ & $1{,}3\cdot 10^{-5}$ \\


\hline


0{,}5 & 114 & $2{,}0\cdot 10^{-5}$ & $6{,}5\cdot 10^{-4}$ &


0{,}04 & 133 & $5{,}0\cdot 10^{-6}$ & $1{,}4\cdot 10^{-5}$ \\


\hline


\end{tabular}


\end{center}


\medskip





Из таблицы видно, что чем меньше шаг дискретизации, тем больше $s$. Причем 


минимальное число итераций $s_0$ достигается


при $h=0{,}1$ , когда $\tau$ принимает некоторое значение,


равное $\tau_0 \approx 0{,}02$.








Следует отметить, что если проверку точности


итерационного процесса проводить только по


невязке $\|r(s)\|\leq \varepsilon$,


то при увеличении $\tau$ $(\tau>50)$ сходимость метода (1.7)


достигается за одну-две итерации, при этом его относительная погрешность


$\|{\overset{s}{\widetilde y}}-y_{\text{точн}}\|/\|y_{\text{точн}}\|$


(т.\,е. погрешность в рамках данной задачи) составляет примерно


0{,}3\% для указанных выше сеточных шагов.








\begin{thebibliography}{99}





\bibitem{1}


Яненко Н.\,Н.


{\it Метод дробных шагов решения многомерных задач математической физики}.


-- Новосибирск: Наука, 1967. -- 195 с.


\bibitem{2}


Самарский А.\,А.


{\it Введение в теорию разностных схем.} -- М.: Наука, 1971. -- 552 с.


\bibitem{3}


Марчук Г.\,И.


{\it Методы вычислительной математики.} -- М.: Наука, 1989. -- 608 с.


\bibitem{4}


Марчук Г.\,И.


{\it Методы расщепления.} -- М.: Наука, 1988. -- 264 с.


\bibitem{5}


Самарский А.\,А., Вабищевич П.\,Н.


{\it Аддитивные схемы для задач математической физики.} --


М.: Наука, 1999. -- 292 с.


\bibitem{6}


Андреев В.\,Б.


{\it Итерационные методы переменных направлений для численного 


решения третьей краевой задачи в $p$-мерном параллелепипеде}


// Журн. вычисл. матем. и матем. физ. -- 1965. --


Т.\,5. -- \No\,4. -- C.\,626--637.


\bibitem{7}


Дьяконов Е.\,Г.


{\it О некоторых итерационных методах решения систем разностных уравнений,


возникающих при решении методом сеток уравнений в частных производных 


эллиптического типа} 


// Вычисл. методы и программирование. -- М. -- 1965. 


-- Вып.\,3. -- C.\,163--194.


\bibitem{8}


Дьяконов Е.\,Г.


{\it Итерационные методы решения разностных аналогов краевых задач 


для уравнений эллиптического типа}


// Киев. Ин-т кибернет. AH УССР. Соврем. числ. методы. -- 1970.


-- Вып.\,4. -- 144 с.


\bibitem{9}


Douglas J., Gunn J.\,E.


{\it A general formulation of alternating direction methods. --- Parabolic


and hyperbolic problems}


// Num. Math. -- 1964. -- V.\,6. -- P.\,428--453.


\bibitem{10}


Ильин В.\,П.


{\it О расщеплении разностных уравнений параболического и 


эллиптического типов}


// Сиб. матем. журн. -- 1965. -- \No\,6. -- C.\,1425--1428.


\bibitem{11}


Birkhoff G., Varga R.\,S.


{\it Implicit alternating direction methods}


// Trans. Amer. Math. Soc. -- 1959. -- V.\,92. -- \No\,1. -- P.\,13--24.


\bibitem{12}


Birkhoff G., Varga R.\,S.


{\it Alternating direction implicit methods}


// Advances in Comp. -- New~York--London:  Academ. Press. -- 1962.


-- V.\,3. -- P.\,189--273.


\bibitem{13}


Белсусан А., Лионс Ж.-Л., Темам Р.


{\it Методы декомпозиции, децентрализации, координации и их приложения}


// Методы вычислительной математики. -- Новосибирск: Наука, 


1975. -- С.\,144--274.


\bibitem{14}


Абрашин В.\,Н., Жадаева Н.\,Г.


{\it Разностные схемы для задач математической физики в области


произвольной формы}


// Дифференц. уравнения и их применение. -- Вильнюс, 1988. -- Вып.\,43. --


С.\,22--30.


\bibitem{15}


Жадаева Н.\,Г.


{\it Многокомпонентный метод переменных направлений решения многомерных


задач для эллиптических уравнений со смешанными производными}


// Дифференц. уравнения. -- 1998. -- Т.\,34. -- \No\,7. --


C.\,948--957.


\bibitem{16}


Абрашин В.\,Н.


{\it Об одном варианте метода переменных направлений решения многомерных 


задач математической физики}


// Дифференц. уравнения. -- 1990. -- Т.\,26. -- \No\,2.


-- C.\,313--323.


\bibitem{17}


Абрашин В.\,Н., Муха В.\,А.


{\it Об одном классе экономичных разностных схем решения многомерных


задач математической физики} 


// Дифференц. уравнения. -- 1992. -- Т.\,28. -- \No\,10. 


-- C.\,1786--1799.


\bibitem{18}


Абрашин В.\,Н.


{\it Об одном итерационном методе решения разностных задач


для эллиптических уравнений}


// Дифференц. уравнения. -- 1998. -- Т.\,31. -- \No\,7.


-- C.\,911--920.


\bibitem{19}


Абрашин В.\,Н.


{\it Многокомпонентные итерационные методы переменных


направлений}


// Матем. моделир. -- 2000. -- Т.\,12. -- \No\,3. -- C.\,45--56.


\bibitem{20}


Самарский А.\,А., Абрашин В.\,Н., Жадаева Н.\,Г.


{\it Аддитивные итерационные методы решения задач математической физики}


// Докл. РАН. -- 2000. -- Т.\,373. -- \No\,6. % -- С.\,734.


\bibitem{21}


Абрашин В.\,Н., Егоров А.\,А., Жадаева Н.\,Г.


{\it О скорости сходимости аддитивных итерационных методов}


// Дифференц. уравнения. -- 2001. -- Т.\,37. -- \No\,7. -- C.\,867--879.


\bibitem{22}


Абрашин В.\,Н., Егоров А.\,А., Жадаева Н.\,Г.


{\it Об одном классе аддитивных итерационных методов}


// Дифференц. уравнения. -- 2001. -- Т.\,37. -- \No\,12. -- C.\,1664--1763.





\end{thebibliography}





\vskip 1cm





\post{\it Институт математики}{\it Поступила}


\post{\it Национальной Академии наук Беларуси}{18.06.2002}


\smallskip


\post{\it Белорусский государственный университет}{}





\label{end}


\end{document}


